\documentclass[memoria_final]{subfiles}
\begin{document}
\edef\@currjob{\jobname}
\def\@mainjob{memoria_final}
\ifx\@currjob\@mainjob
  \clearpage
  \part{Fase 3: Controlador y Sistema Completo}
\else
  \hypertarget{inicio}{}

  \begin{titlepage}
    \thispagestyle{empty}
    \begin{tikzpicture}[remember picture,overlay]
      \fill[white](current page.south west)rectangle(current page.north east);
      \fill[MainBlue]
        (current page.north west)rectangle
        ([yshift=-0.22\paperheight]current page.north east);
      \draw[MainBlue!40,line width=2pt]
        ([yshift=-0.22\paperheight]current page.north west)--
        ([yshift=-0.22\paperheight]current page.north east);
    \end{tikzpicture}

    \vspace*{-1cm}
    \begin{flushleft}
      {\fontsize{9}{12}\selectfont\color{white!80}
        GRADO EN INGENIERÍA ELECTRÓNICA DE COMUNICACIONES}\\[1ex]
      {\fontsize{24}{28}\selectfont\bfseries\color{white}
        Diseño Digital 2}\\[1ex]
      {\fontsize{11}{14}\selectfont\color{white!80}
        Curso académico 2025--2026}
    \end{flushleft}

    \vspace*{0.11\paperheight}
    \begin{center}
      {\Large\color{MainBlue!80}\textbf{BLOQUE II}
        \quad\textbullet\quad Controlador y Sistema Completo}\\[0.8cm]

      {\fontsize{30}{38}\selectfont\bfseries\color{MainBlue}
        Memoria del Proyecto}\\[0.2cm]
      {\fontsize{22}{28}\selectfont\bfseries\color{MainBlue!80}
        Bloque II --- Fase 3}\\[0.5cm]

      {\fontsize{12}{18}\selectfont\color{black!60}
        Diseño, implementación y verificación en VHDL-93}\\[0.2cm]
      {\fontsize{11}{14}\selectfont\ttfamily\color{black!50}
        controlador \quad calculadora \quad tb\_calculadora}\\[0.4cm]

      {\fontsize{11}{14}\selectfont\bfseries\color{MainBlue}
        Versión 1.0 --- Mayo 2026}

      \vspace{1cm}
      \renewcommand{\arraystretch}{1.5}
      \begin{tabular}{@{}r@{\hspace{1em}}l@{}}
        {\small\color{MainBlue}\faLayerGroup}\enspace\textbf{Fase:}
          & Bloque II --- Fase 3 \\
        {\small\color{MainBlue}\faCalendar}\enspace\textbf{Entrega:}
          & 28 de Mayo de 2026 \\
        {\small\color{MainBlue}\faMicrochip}\enspace\textbf{Plataforma:}
          & DECA MAX10 + conector XDECA \\
        {\small\color{MainBlue}\faClock}\enspace\textbf{Reloj:}
          & 50\,MHz \\
      \end{tabular}

      \vspace{1cm}
      \begin{tcolorbox}[
        enhanced,width=0.80\textwidth,
        colback=SoftBlue,colframe=MainBlue,
        boxrule=0.8pt,arc=4pt,
        title={\centering\textbf{\faUsers\quad Miembros del Grupo}},
        coltitle=white,fonttitle=\large]
        \begin{center}
          \renewcommand{\arraystretch}{1.3}
          \begin{tabular}{@{}ll@{}}
            \textbf{Lucía Llana Carmona}   &\small\texttt{lucia.llana@alumnos.upm.es}\\
            \textbf{María Moya Marcilla}    &\small\texttt{maria.moya.marcilla@alumnos.upm.es}\\
            \textbf{Pedro Gil Bolaños}      &\small\texttt{pedro.gilb@alumnos.upm.es}\\
            \textbf{Pedro Osma Moya}        &\small\texttt{p.osma@alumnos.upm.es}\\
            \textbf{Eduardo Muñoz de Maya}  &\small\texttt{eduardo.mdemaya@alumnos.upm.es}\\
          \end{tabular}
        \end{center}
      \end{tcolorbox}
    \end{center}
    \vfill
    \begin{center}
      {\fontsize{8}{10}\selectfont\color{black!40}
        Grado en Ingeniería Electrónica de Comunicaciones
        \quad$\cdot$\quad Diseño Digital 2}
    \end{center}
    \vspace{0.5cm}
  \end{titlepage}

  \newpage
  \tableofcontents
  \newpage
\fi% fin \ifx\@currjob\@mainjob
\fancyhead[R]{\small\textbf{Memoria --- Fase 3 (Controlador y Sistema Completo)}}

% ===================================================================
\section{Introducción y objetivos}
% ===================================================================

En esta tercera fase se diseña el módulo \textbf{controlador} (la máquina de
estados principal de la calculadora) y el módulo \textbf{calculadora} (top-level
estructural que integra todos los bloques de las tres fases). Finalmente se
verifica el sistema completo mediante el testbench \texttt{tb\_calculadora}.

% ===================================================================
\section{Controlador (\texttt{controlador.vhd})}
% ===================================================================

\subsection{Máquina de estados}

El controlador implementa una FSM Moore de 4 estados:
\begin{itemize}
  \item \textbf{OP1}: espera el primer operando (teclas 0-9, signo C,
    operador A/D/E).
  \item \textbf{OP2}: espera el segundo operando.
  \item \textbf{WAIT\_BCD}: espera a que \texttt{bin\_to\_bcd} termine la
    conversión (\texttt{done = '1'}).
  \item \textbf{RES}: muestra el resultado en los displays.
\end{itemize}

% ===================================================================
\section{Integración: \texttt{calculadora.vhd}}
% ===================================================================

El módulo top-level \texttt{calculadora} conecta estructuralmente todos los
bloques en el camino de datos completo: teclado $\to$ controlador $\to$
BCD2BIN $\to$ ALU $\to$ BIN2BCD $\to$ displays.

% ===================================================================
\section{Verificación: \texttt{tb\_calculadora.vhd}}
% ===================================================================

El testbench prueba 4 operaciones del sistema completo:
\begin{itemize}
  \item $123 + 456 = 579$
  \item $999 - 1 = 998$
  \item $12 \times 12 = 144$
  \item $(-50) - 200 = -250$
\end{itemize}

% ===================================================================
\section{Conclusiones}
% ===================================================================

\begin{tcolorbox}[colback=SoftBlue,colframe=MainBlue,boxrule=0.8pt,arc=2pt,
  left=0.5cm,right=0.5cm,top=0.3cm,bottom=0.3cm]
  \small\faInfoCircle\enspace Esta sección se completará tras la entrega
  de la Fase 3 (28 de Mayo de 2026).
\end{tcolorbox}

\end{document}
